\scanpage{83}
est muni d'un élément \inserted{$\handwritten{\underline}{f}_q$} de $G_{\deleted{a}\handwritten{\unsure{\nu|a}}}(K_{\deleted{a}\handwritten{\nu}})$\inserted{(\handwritten{``}Frobenius relatif à $\unsure{\handwritten{\underline}{F}|\handwritten{\underline}{\mathbb{F}}}_q$'')}. Ce dernier n'est autre d'ailleurs que $\deleted{\handwritten{\underline}{f}_q}$ $\handwritten{\underline}{f}^{\deleted{a}\handwritten{\nu}}$, où $\handwritten{\underline}{f}$ est le générateur canonique de $G$ (le Frobenius absolu \handwritten{!}). \deleted{\unsure{O|C} on voit donc que $G_a$ est encore engendré par} Notons que $a\unsure{\operatorname{Id}|\operatorname{id}}_G$ est un monomorphisme (car $x\handwritten{\mapsto}x^{\deleted{a}\handwritten{\nu}}$ est un épimorphisme dans $\handwritten{\overline}{\handwritten{\underline}{Q}}_{\handwritten{p}}^{\handwritten{*}}$ ), donc induit un isomorphisme entre $G_{K_{\deleted{a}\handwritten{\nu}}}$ et $G_{\deleted{a}\handwritten{\nu}}\handwritten{\simeq}G(\handwritten{\underline}{F}_q)$,transformant le frobenius \deleted{\unsure{relati|relatif}} absolu en le frobenius relatif à $\handwritten{\underline}{F}_q$. Grâce à ceci, on voit :

\handwritten{\uline}{Proposition} 10.1. La catégorie \inserted{tann.} $\operatorname{Fcriso}(\handwritten{\underline}{F}_q)\handwritten{\otimes}_{\handwritten{\underline}{Q}_p}K_a$ est équivalente à la catégorie \inserted{tann.} $\operatorname{Fcriso}(\handwritten{\underline}{F}_p)\handwritten{\otimes}_{\handwritten{\underline}{Q}_p}K_a$,\deleted{\handwritten{\unsure{(|/}}} \margin{\handwritten{(i.e. la catégorie des couples $(E,F'_E)$ d'un vectoriel sur $K_\nu$ et d'un autom. de ce vectoriel....)}}
avec $\handwritten{\underline}{f}_q$ correspondant à $\handwritten{\underline}{f}^{\deleted{a}\handwritten{\nu}}$ (où on peut, si on veut, considérer $\handwritten{\underline}{f}_q$ et $\handwritten{\underline}{f}$ comme des endomorphismes des foncteurs identiques de $\operatorname{Fcriso}(\handwritten{\underline}{F}_q)$ resp. $\operatorname{Fcriso}(\handwritten{\underline}{F}_p)$)), et correspondance des foncteurs fibres canoniques. Le foncteur déduit du changement de base $\handwritten{\underline}{F}_p\handwritten{\to}\handwritten{\underline}{F}_q$

\[
\operatorname{Fcriso}(\handwritten{\underline}{F}_p)\handwritten{\otimes}_{\handwritten{\underline}{Q}_p}K_{\deleted{a}\handwritten{\nu}}\handwritten{\longrightarrow}\operatorname{Fcriso}(\handwritten{\underline}{F}_q)\handwritten{\otimes}_{\handwritten{\underline}{Q}_p}K_{\deleted{a}\handwritten{\nu}}
\]
s'identifie, par l'équivalence précédente, au foncteur qui à tout couple $(E,F_E)$ d'un vectoriel sur $K_{\deleted{a}\handwritten{\nu}}$ et d'un automorphisme dudit, associe le couple $(E,F_E{}^{\deleted{a}\handwritten{\nu}})$\handwritten{;}
\inserted{\handwritten{analogue pour $\operatorname{Fcriso}(F_{\unsure{q|1}})-\operatorname{Fcriso}(F_q)\otimes_{\mathbb{Q}_p}K_\nu$, il s'identifie : $(E,F_E)\longmapsto(E,F_E{}^\nu)$.}}

10.2. Comme on connait $\handwritten{\underline}{G}(\handwritten{\underline}{F}_q)$ comme une sous-gerbe $\handwritten{\underline}{G}_{\deleted{a}\handwritten{\nu}}$ de $\handwritten{\underline}{G}=\handwritten{\underline}{G}(\handwritten{\underline}{F}_p)$, et que celle-ci contient $\handwritten{\underline}{H}$, on connait donc en principe l'inclusion
\[
\handwritten{\underline}{H}\handwritten{\longrightarrow}\handwritten{\underline}{G}_{\deleted{a}\handwritten{\nu}}=\handwritten{\underline}{G}(\handwritten{\underline}{F}_q)\quad ,
\]
donc aussi le foncteur
\[
\operatorname{Rep}(\handwritten{\underline}{G}_{\deleted{a}\handwritten{\nu}})\handwritten{\simeq}\operatorname{Fcriso}(\handwritten{\underline}{F}_q)\handwritten{\longrightarrow}\operatorname{Rep}(\handwritten{\underline}{H})\quad ,
\]
\deleted{\inserted{et} Considérons le foncteur correspondant}
\[
\deleted{\operatorname{Fcriso}(\handwritten{\underline}{F}_q)\handwritten{\otimes}_{\handwritten{\underline}{Q}_p}K_{\underset{\deleted{\unsure{a|x}}}{a}}\quad\operatorname{Rep}(\handwritten{\underline}{H})\handwritten{\otimes}_{\handwritten{\underline}{Q}_p}K_a\ .}
\]
L'effet de ce\deleted{s} foncteur\deleted{s} sur les classes d'isomorphie d'éléments est connu (comme chaque fois qu'il s'agit d'un morphisme d'une \deleted{gerbe de} \unsure{Catégorie|catégorie|Oatégorie} tannakienne à lien commutatif dans une autre à lien de type multiplicatif) par la connaissance de l'homomorphisme de liens correspondant
% archive sheet number: 82
